\scanpage{4}
\uline{Dém} On sait déjà (\unsure{réf|bof}) que $J$ est
q. coh. et $J\vert{}X'=J'$, donc $G$ est q. \unsure{coh|ch}
et $G|X'=G'$. Prouvons que $G$ est $V$-plat,
i.e. que l'uniformisante $\pi$ de $V$ est
non diviseur de $0$ dans $G$. \deleted{\unsure{Sinon,}}
Soit en effet \deleted{\unsure{$J$} \unsure{$K$}} $\unsure{\mathcal N|N}$ \deleted{\unsure{l'élément}} noyau
de $\pi$ dans $G$, alors $G/\mathcal N$ est un
\deleted{faisceau}\inserted{Module} quotient \deleted{\unsure{qu}} de $G$ qui induit
encore $G'$ sur $X'$, donc d'après la
condition minimale de $\deleted{\unsure{\mathcal N}}G$, \deleted{\illegible{1}} on a $\mathcal N=0$,
ce qui prouve (i).

Pour (ii), on note que si
$G_1$ est le plus petit \unsure{module|Module} quotient
de $G$ induisant $G'$ sur $X'$, \deleted{\unsure{Alors}}
de sorte que $G_1$ est \inserted{un quotient de $F$} majoré par $G$,
\margin{$G\to G_1$}alors le noyau de \inserted{\unsure{l'épim|libérer}} $G\to G_1$ est \inserted{\unsure{encore|aussi|inclus}} plat sur $V$,
et induit $0$ sur $X'$, il suffit de
prouver qu'un tel Module q. coh. est
Nul. \unsure{Cela|Or} est immédiat, puisque
\[
N\otimes_V K=0.\qquad \text{Cqfd.}
\]
% archive sheet number: 3

